\section{Timing And Further Reading}\label{sec:further-reading}

Spectre prepares each new spectrum over its 1024-sample update interval,
about $21.3\,\mathrm{ms}$ at $48\,\mathrm{kHz}$. Each spectrum already covers
about $42.7\,\mathrm{ms}$ of audio at that rate. Average and Rack's display
refresh add further visible delay.

Analysis settings take effect at the start of the next frame and affect new
columns. Turning Run off pauses any unfinished frame; resuming completes it
before starting another. Changing Rack's sample rate discards unfinished
analysis. Allow new history to replace older columns before comparing results
across that change.

Analysis runs on Rack's audio engine thread. Work is spread across sample
calls to reduce processing bursts, but the module still contributes to the
patch's CPU load. This does not guarantee that a heavily loaded patch will
meet its audio deadlines.

For the mathematical background, FFT algorithms, scheduling proofs, and
performance experiments, see the accompanying
\href{https://github.com/Kautenja/ArhythmeticUnits-Fourier/tree/main/docs/whitepaper}{\textit{Whole-Pipeline Scheduling for Real-Time Spectral Analysis}}.
Its implementation appendix gives the supporting algorithms; this manual
focuses on operating the module.
